\begin{table}[htbp]\centering
\begin{threeparttable}
\caption{Names given to the two classes of divergent sequences in the three systems, compared with the main type of their culture.}
\label{tab:divergent}
\small
\begin{tabular}{p{0.50\linewidth}p{0.18\linewidth}p{0.18\linewidth}}
\toprule
 & Short 5.8S & Same length\\
\midrule
Sequences (cultures) & 43 (24) & 8 (5)\\
VT: same as the culture's main type / other / no name & 35 / 7 / 1 & 8 / 0 / 0\\
LSU: same species / another species / no name & 5 / 1 / 31 & 2 / 5 / 1\\
LSU best-hit identity, median (main type 98.9\%) & 94.9\% & 97.2\%\\
SH: same / other SH / no name & 0 / 28 / 7 & 0 / 4 / 1\\
SH best-hit identity, median (main type 99.1\%) & 99.8\% & 98.6\%\\
\bottomrule
\end{tabular}
\begin{tablenotes}\footnotesize
\item Each divergent sequence compared with the most common name of its culture's main-type sequences (culture
reference version~4). Sequences in cultures without a main-type name in a system are left out of that
system's row (short 5.8S: VT 0, LSU 6, SH 8;
same length: VT 0, LSU 0, SH 3).
\end{tablenotes}
\end{threeparttable}
\end{table}
